% Algorithm description for NSGA-III nutrition optimization
\section{NSGA-III Algorithm Details}

\subsection{Pseudocode}

\begin{algorithm}[H]
\caption{NSGA-III for Nutrition Optimization}
\begin{algorithmic}[1]
\STATE Initialize population $P_0$ with $N$ random solutions
\STATE Generate reference directions $\mathbf{R}$ using Das-Dennis method
\FOR{$t = 1$ to $T$}
    \STATE Create offspring population $Q_t$ using genetic operations
    \STATE Combine parent and offspring: $R_t = P_t \cup Q_t$
    \STATE Perform non-dominated sorting on $R_t$
    \STATE Initialize next population: $P_{t+1} = \emptyset$
    \STATE $i = 1$
    \WHILE{$|P_{t+1}| + |F_i| \leq N$}
        \STATE Add front $F_i$ to $P_{t+1}$
        \STATE $i = i + 1$
    \ENDWHILE
    \IF{$|P_{t+1}| < N$}
        \STATE Associate solutions in $F_i$ with reference directions
        \STATE Select remaining solutions using niche preservation
    \ENDIF
\ENDFOR
\RETURN Final population $P_T$
\end{algorithmic}
\end{algorithm}

\subsection{Reference Direction Generation}

The Das-Dennis systematic approach generates reference directions as follows:

\begin{algorithm}[H]
\caption{Das-Dennis Reference Direction Generation}
\begin{algorithmic}[1]
\STATE Input: Number of objectives $M$, number of partitions $H$
\STATE Initialize reference set $\mathbf{R} = \emptyset$
\FOR{all combinations of $w_1, w_2, \ldots, w_M$}
    \IF{$\sum_{i=1}^{M} w_i = 1$ and $w_i \in \{0, \frac{1}{H}, \frac{2}{H}, \ldots, 1\}$}
        \STATE $\mathbf{r} = \frac{\mathbf{w}}{\|\mathbf{w}\|}$
        \STATE Add $\mathbf{r}$ to $\mathbf{R}$
    \ENDIF
\ENDFOR
\RETURN Reference directions $\mathbf{R}$
\end{algorithmic}
\end{algorithm}

\subsection{Solution Association}

Each solution is associated with the nearest reference direction:

\begin{equation}
d(\mathbf{z}, \mathbf{r}) = \arccos\left(\frac{\mathbf{z}^T \mathbf{r}}{\|\mathbf{z}\| \|\mathbf{r}\|}\right)
\end{equation}

where $\mathbf{z}$ is the normalized objective vector and $\mathbf{r}$ is a reference direction.

\subsection{Niche Preservation}

When multiple solutions are associated with the same reference direction, niche preservation selects the solution with the smallest perpendicular distance:

\begin{equation}
d_p(\mathbf{z}, \mathbf{r}) = \|\mathbf{z} - \frac{\mathbf{z}^T \mathbf{r}}{\|\mathbf{r}\|^2} \mathbf{r}\|
\end{equation}

\subsection{Genetic Operations}

\subsubsection{Crossover}

We use simulated binary crossover (SBX) with distribution index $\eta_c = 20$:

\begin{equation}
\beta = \begin{cases}
(2u)^{\frac{1}{\eta_c + 1}} & \text{if } u \leq 0.5 \\
\left(\frac{1}{2(1-u)}\right)^{\frac{1}{\eta_c + 1}} & \text{otherwise}
\end{cases}
\end{equation}

\subsubsection{Mutation}

Polynomial mutation with distribution index $\eta_m = 20$:

\begin{equation}
\delta = \begin{cases}
(2r)^{\frac{1}{\eta_m + 1}} - 1 & \text{if } r < 0.5 \\
1 - (2(1-r))^{\frac{1}{\eta_m + 1}} & \text{otherwise}
\end{cases}
\end{equation}

\subsection{Computational Complexity}

The overall computational complexity of NSGA-III is:

\begin{itemize}
    \item \textbf{Non-dominated sorting}: $O(MN^2)$ per generation
    \item \textbf{Reference direction association}: $O(MN|\mathbf{R}|)$ per generation
    \item \textbf{Total complexity}: $O(T \cdot M \cdot N^2)$ for $T$ generations
\end{itemize}

where $M$ is the number of objectives, $N$ is the population size, and $|\mathbf{R}|$ is the number of reference directions.

\subsection{Parameter Settings}

\begin{table}[H]
\centering
\caption{NSGA-III Parameter Settings}
\label{tab:nsga3_params}
\begin{tabular}{lc}
\toprule
Parameter & Value \\
\midrule
Population size & 100 \\
Number of generations & 200 \\
Crossover probability & 0.9 \\
Mutation probability & 0.1 \\
Distribution index (crossover) & 20 \\
Distribution index (mutation) & 20 \\
Reference direction partitions & 12 \\
\bottomrule
\end{tabular}
\end{table}

\subsection{Convergence Analysis}

The convergence behavior of NSGA-III can be analyzed through:

\begin{enumerate}
    \item \textbf{Hypervolume indicator}: Measures the volume of objective space dominated by the Pareto front
    \item \textbf{Generational distance}: Average distance from solutions to the true Pareto front
    \item \textbf{Spread indicator}: Measures the distribution of solutions along the Pareto front
\end{enumerate}

For our nutrition optimization problem, we observe:

\begin{itemize}
    \item Rapid convergence in early generations (1-50)
    \item Steady improvement in middle generations (51-150)
    \item Fine-tuning in later generations (151-200)
\end{itemize}

\subsection{Solution Diversity}

NSGA-III maintains solution diversity through:

\begin{enumerate}
    \item \textbf{Reference point distribution}: Ensures coverage across the objective space
    \item \textbf{Niche preservation}: Prevents clustering around popular reference directions
    \item \textbf{Environmental selection}: Balances convergence and diversity
\end{enumerate}

The diversity metric is measured as:

\begin{equation}
Diversity = \frac{1}{N(N-1)} \sum_{i=1}^{N} \sum_{j=i+1}^{N} \|\mathbf{x}_i - \mathbf{x}_j\|
\end{equation}

where $\mathbf{x}_i$ and $\mathbf{x}_j$ are solution vectors in the decision space. 